\documentclass[11pt,a4paper]{article}

\usepackage[utf8]{inputenc}
\usepackage[T1]{fontenc}
% lmodern unavailable in this TeX install; disable microtype font expansion instead
\usepackage[margin=1in]{geometry}
\usepackage{amsmath,amssymb,amsthm}
\usepackage{booktabs}
\usepackage{listings}
\usepackage{xcolor}
\usepackage[protrusion=true,expansion=false]{microtype}
\usepackage[hidelinks]{hyperref}

\lstset{
  basicstyle=\ttfamily\footnotesize,
  keywordstyle=\color{blue!60!black},
  commentstyle=\color{green!40!black}\itshape,
  breaklines=true,
  frame=single,
  framesep=4pt,
  showstringspaces=false,
  columns=fullflexible,
}

\theoremstyle{plain}
\newtheorem{proposition}{Proposition}
\newtheorem{lemma}{Lemma}
\theoremstyle{definition}
\newtheorem{definition}{Definition}
\theoremstyle{remark}
\newtheorem*{remark}{Remark}

\newcommand{\FF}{\mathbb{F}}

\title{\textbf{Concentric Cubes Are Plumbing, Not Security}\\[0.4em]
\large Design and Cryptanalysis of \texttt{cube3d-sponge}, a Three-Dimensional\\
Cellular-Automaton Sponge Permutation}

\author{Wofl \\ \small whispr.dev \\ \small \texttt{<contact>}}

\date{\today}

\begin{document}
\maketitle

\begin{abstract}
A recurring intuition in the cellular-automaton (CA) cryptography literature holds
that adding spatial dimensions, neighbourhood radius, or nested structure to a CA
makes the resulting cipher ``more complex and hence more secure.'' We take that
intuition seriously enough to test it. Starting from a proposal to fold plaintext
into a three-dimensional grid of concentric cubic shells and evolve a CA radiating
outward from the centre, we show that the geometric intuition is \emph{orthogonal}
to security: apparent complexity contributes nothing that is not already captured by
nonlinearity, invertibility, symmetry-breaking, and diffusion uniformity. We then
show that the concentric-shell geometry \emph{does} have one legitimate cryptographic
role --- as a sponge rate/capacity boundary --- and build \texttt{cube3d-sponge}, an
$8\times8\times8$ byte-state ($4096$-bit) permutation combining a checkerboard
generalised Feistel network over a 3-torus, a $\chi$-family nonlinear map, a per-cell
byte rotation, and a coordinate shear, wrapped in a sponge whose rate is the outer
cubic shell and whose capacity is the inner shells.

Our contribution is primarily \emph{negative and methodological}. We report a
complete cryptanalytic ladder against the construction, obtained from four
independent harnesses and cross-validated on two operating systems: key recovery by
cube attack succeeds only to round~2 of~12; deterministic zero-sum distinguishers
derived from provable degree upper bounds reach round~4 with practical cubes
($2^{16}$) and round~5 with a $2^{64}$ cube; statistical differential bias is
detectable to round~6. We document two design bugs the harnesses caught that would
otherwise have produced confident nonsense --- a checkerboard 2-colouring that is
inconsistent on an odd-length torus, and bitwise operations that left the eight
bit-planes of the state \emph{provably non-interacting}, capping avalanche at exactly
$512/2/4096 = 6.25\%$ --- and two harness bugs that produced impossible
(non-monotonic) measurements before being fixed. We argue that the discipline of
building the falsifier before trusting the design is the transferable result here,
and we explicitly decline to claim that \texttt{cube3d-sponge} is secure. It is
merely \emph{not obviously broken}, which is a far weaker statement, and we recommend
Ascon or ChaCha20-Poly1305 for any real use.
\end{abstract}

\tableofcontents

\section{Introduction}

\subsection{The question}

Cellular automata are seductive to cryptographers. They are local, parallel, cheap in
hardware, and they generate visually chaotic output from trivially simple rules. It
is tempting --- and the literature shows it has been tempting for forty years --- to
read that visual chaos as cryptographic strength.

This work began with a concrete proposal in that tradition: given prior work
analysing data plotted as two-dimensional grids under CA evolution, would folding
data into a \emph{three-dimensional} grid --- a central cube surrounded by concentric
shells of smaller cubes, with the CA radiating outward in all directions --- make the
technique ``way more complex and hence more secure''?

The question deserves a precise answer, because the reasoning inside it is the single
most common failure mode in the CA- and chaos-cryptography literature. The answer has
three parts:

\begin{enumerate}
  \item \textbf{No, not as stated.} Complexity is not security. Security is the
    analysed work factor for an adversary who knows the entire construction and lacks
    only the key (Kerckhoffs~\cite{kerckhoffs}). Adding dimensions, shells, or radius
    adds state and cost; it does not add any of the properties that actually resist
    cryptanalysis. A larger linear system is still solved by Gaussian elimination.
  \item \textbf{But the shape is right.} A 3D local nonlinear permutation iterated
    with round constants is, almost exactly, Keccak~\cite{keccak}. Keccak's state
    \emph{is} a $5\times5\times64$ cube of bits. The difference between Keccak and a
    naive radiating CA is a small, enumerable set of properties --- not the geometry.
  \item \textbf{And one piece of the geometry earns its keep.} Concentric shells are
    a natural \emph{rate/capacity boundary}: outer shell public, inner shells never
    directly output. That is the sponge construction~\cite{spongeindiff}, and it is a
    legitimate, analysed structural role for ``layers.''
\end{enumerate}

\subsection{What we did}

We built the construction that (3) suggests, and then we attacked it. The point was
never to produce a competitive cipher --- we say plainly in
Section~\ref{sec:limitations} that nobody should use this --- but to find out what a
disciplined cryptanalytic process reports about a design whose provenance is
``geometric intuition, corrected.''

Specifically, we contribute:

\begin{itemize}
  \item \textbf{A design} (Section~\ref{sec:design}) that isolates the geometric idea
    from the security-relevant machinery: concentric shells serve \emph{only} as the
    rate/capacity split, while diffusion uniformity is handled independently by
    toroidal wrapping. Conflating those two roles --- which the original ``radiate
    from the centre'' framing does --- makes shell boundaries simultaneously diffusion
    weak points and security boundaries.
  \item \textbf{A cryptanalytic ladder} (Section~\ref{sec:results}) from four
    independent harnesses: avalanche/diffusion, statistical differential and linear
    distinguishers, a cube attack with real key recovery, and provable algebraic
    degree bounds yielding deterministic zero-sums.
  \item \textbf{A record of the bugs} (Sections~\ref{sec:bugs} and~\ref{sec:harnessbugs}),
    both in the design and in the harnesses. We consider this the most transferable
    part of the paper. Every one of the four was caught by a measurement that produced
    a \emph{physically impossible} number, not by inspection.
  \item \textbf{An explicit refusal to overclaim.} Upper bounds on degree prove that
    distinguishers exist; they never prove that they do not. Negative attack results
    mean ``this attack, at this budget, found nothing.''
\end{itemize}

\subsection{A note on a corrected novelty claim}
\label{sec:novelty}

Early in this work we believed the concentric-cube geometry to be structurally novel.
A literature search conducted for this paper shows that belief was too strong, and we
correct it here rather than quietly dropping it. Martín del Rey~\cite{rey2015}
encrypts 3D solid objects using a three-dimensional chaotic Cat map for confusion
followed by a \emph{second-order memory reversible 3D cellular automaton} for
diffusion, iterated over several rounds --- which is, almost point for point, the
architecture we independently proposed as the ``principled path.'' 3D CA image and
object encryption is an established sub-literature~\cite{rey2015,chen2004,3dca2016,carev2026}.
What remains arguably new in our construction is narrower: the checkerboard Feistel
over a 3-torus as an invertibility mechanism, and the use of concentric shells
specifically as a sponge rate/capacity boundary rather than as a diffusion structure.
We make no strong novelty claim beyond that, and we note that novelty of geometry
would not, in any case, imply novelty of security argument.

\section{Related work}
\label{sec:related}

\subsection{Cellular-automaton cryptography and its cryptanalytic record}

Wolfram proposed elementary CA rule~30 as the basis of a Vernam-like stream cipher at
CRYPTO~1985~\cite{wolfram85}, motivated by its Class~III chaotic dynamics: the key is
the initial configuration of a ring of cells, and the sequence of states of a
distinguished cell is the keystream. Meier and Staffelbach broke it at
EUROCRYPT~1991~\cite{meier91} with a correlation attack that reconstructs the seed
despite the sequence passing classical statistical randomness tests. This is the
canonical lesson of the field, and it is worth stating in its strongest form:
\emph{passing statistical tests is not evidence of cryptographic strength}. Rule~30
passed; rule~30 fell.

Subsequent work has attempted to repair the approach --- by enlarging the
neighbourhood radius~\cite{4nbr}, by using maximum-period nonlinear
CA~\cite{maiti17}, by hybridising linear and nonlinear rules --- with mixed and often
short-lived success. Mariot and Leporati's recent survey~\cite{camariot} of a decade
of CA-based cryptography is candid about the field's methodological problems, noting
in particular that most work relies on ad~hoc arguments specific to individual rules
that generalise poorly.

It is worth recording the one unambiguous success story: the $\chi$ mapping, a
degree-2 shift-invariant nonlinear CA-like transformation introduced by Daemen et
al.~\cite{daemen94}, is the nonlinear core of Keccak~\cite{keccak} and, in affine
equivalent form, of Ascon~\cite{asconspec}. CA ideas did produce a
standards-grade primitive --- but as \emph{one layer} inside a carefully analysed
round function, never as the whole cipher.

\subsection{The chaos/CA image-encryption literature}

The proximate context for the original question is the large literature on
chaos- and CA-based image encryption. Its architectural template is due to
Fridrich~\cite{fridrich98}: a permutation (confusion) stage scrambles pixel
\emph{positions} using a 2D or 3D chaotic map, followed by a diffusion stage that
modifies pixel \emph{values}, with the pair iterated for several rounds. This
directly instantiates Shannon's confusion/diffusion principle~\cite{shannon49}.

The cryptanalytic record of this literature is poor, and it is important to state so
plainly because sheer publication volume can be mistaken for validation. Fridrich's
own founding scheme was cryptanalysed by Solak et al.~\cite{solak10}, with an attack
that extends to structurally similar designs including the widely-cited 3D chaotic
Cat map scheme of Chen et al.~\cite{chen2004}. Li and collaborators have broken a
long series of such schemes~\cite{li09compound,li11breaking,li18ieaie,limodified},
frequently with a handful of chosen plaintexts. Most damningly, Li, Chen and
Lo~\cite{lipermonly} give a \emph{generic} result: all permutation-only multimedia
ciphers are practically insecure against known/chosen-plaintext attack, requiring only
$O(\log_L(MN))$ plaintexts to recover at least half of a plaintext of size $M\times N$
with $L$ value levels.

Two conclusions follow, both directly relevant to the question this paper answers.
First, ``cascading more chaotic maps / more CA layers'' does not compose weak
primitives into a strong one. Second --- and this is the specific trap in the original
proposal --- a \emph{fixed, public} folding of data into a grid is a known permutation,
and a known permutation contributes exactly zero entropy.

\subsection{Three-dimensional CA cryptography}

Prior 3D CA cryptography exists and is not marginal. Beyond
Rey~\cite{rey2015} discussed in Section~\ref{sec:novelty}, there is work rendering
images into 3D lattices with chaotic Cat map confusion and 3D CA
diffusion~\cite{3dca2016}, 3D image encryption via integral imaging and CA
transforms, and a 2026 review covering 1D, 2D and 3D CA image encryption
comprehensively~\cite{carev2026}. The recurring pattern is that the third dimension is
used as \emph{extra scrambling volume}. We are not aware of prior work using
concentric shells of a 3D CA state as a sponge rate/capacity boundary, which is the
one structural element we would defend as new --- while noting again that structural
novelty is not a security argument.

\subsection{Sponge constructions}

The sponge construction~\cite{spongeindiff} maintains a state split into a
\emph{rate} $r$ (absorbed into and squeezed from) and a \emph{capacity} $c$ (never
touched directly by the outside world), with $b = r + c$. Security against generic
attacks is governed by $c$, not $b$. Keccak~\cite{keccak}, standardised as
SHA-3~\cite{fips202}, instantiates this with a $b = 1600$-bit permutation whose state
is a $5\times5\times64$ array of bits, and whose round is
$\iota\circ\chi\circ\pi\circ\rho\circ\theta$. Ascon~\cite{asconspec}, selected by NIST
for lightweight cryptography and specified in SP~800-232~\cite{sp800232}, uses a
320-bit sponge/duplex with a 5-bit S-box that is an affine equivalent of Keccak's
$\chi$, applying 12 rounds for initialisation/finalisation and 8 for the data path.

The relevance to the original question is direct. Keccak is a three-dimensional,
local, nonlinear, iterated permutation --- the intuition ``operate a 3D local rule
repeatedly on a cube of data'' is not wrong, it is \emph{already standardised}. What
separates Keccak from a naive radiating CA is a short list: a nonlinear step ($\chi$),
a symmetry-breaking round constant ($\iota$), a bit-transposition step ($\rho$, $\pi$)
and a linear diffusion step ($\theta$) --- plus, decisively, a decade of adversarial
analysis.

\subsection{Cryptanalytic tools}

\paragraph{Cube attacks.} Dinur and Shamir~\cite{cube09} observe that any output bit
of a keyed primitive is a polynomial over $\FF_2$ in key and public (IV/nonce)
variables. Summing the output over all $2^d$ assignments of a chosen $d$-subset $I$ of
public variables (the \emph{cube}) annihilates every monomial not divisible by the
cube monomial $t_I$, leaving the \emph{superpoly} $p_{S(I)}$. If the total degree is
at most $d+1$, the superpoly is linear in the key bits, yielding a free linear
equation. Enough independent linear superpolys recover the key by Gaussian
elimination. Linearity is tested probabilistically by the Blum--Luby--Rubinfeld
test~\cite{blr}: a linear $p$ satisfies $p(x)\oplus p(y)\oplus p(0) = p(x\oplus y)$.
Cube testers~\cite{cubetesters} extend the idea to distinguishers.

\paragraph{Division property.} Todo~\cite{todo15} introduced the division property as
a generalised integral property; Todo and Morii~\cite{todo16} refined it to the
bit-based division property (BDP). Xiang et al.~\cite{xiang16} made it practical at
realistic block sizes by modelling division-trail propagation as Mixed Integer Linear
Programming, solved with off-the-shelf solvers. Division property gives the current
state of the art in \emph{provable} degree upper bounds and zero-sum
distinguishers, and Hu et al.'s monomial-prediction
formulation~\cite{hu20} makes degree evaluation exact.

\paragraph{Numeric mapping.} Liu~\cite{liu17} introduced numeric mapping, a cheap
degree-propagation technique for NFSR-based cryptosystems. It propagates per-bit
degree upper bounds through $\deg(a\oplus b)\le\max(\deg a,\deg b)$ and
$\deg(a\wedge b)\le \deg a + \deg b$. It is sound but loose. Zhang et
al.~\cite{zhang21} compare Boura--Canteaut's formula~\cite{boura11}, Carlet's formula,
Liu's numeric mapping and Todo's division property, proving division property is never
worse than the first two. Hu et al.~\cite{hu20} quantify the gap dramatically: the
exact degree of 834-round Trivium is 78, whereas Liu's numeric mapping bound already
reaches 79 at 793 rounds. We use numeric mapping in this work, and
Section~\ref{sec:limitations} is explicit that this is a tractability compromise, not
a choice of the best tool.

\paragraph{Degree bounds.} Canteaut and Videau~\cite{canteaut02} bound the degree of a
composition; Boura, Canteaut and De Cannière~\cite{boura11} give bounds specialised to
Keccak-like permutations that are tighter than the naive $(\deg F)^R$ near saturation.

\section{Design of \texttt{cube3d-sponge}}
\label{sec:design}

\subsection{State geometry}

The state is an $8\times8\times8$ array of bytes: $512$ cells, $4096$ bits. Cell
$(x,y,z)$ has index $i = x + 8y + 64z$.

\paragraph{Why $N=8$ and not $7$.} Our first implementation used $N=7$ to obtain a
single centre cell, matching the ``central cube'' of the original proposal literally.
This is \emph{mathematically broken}, and the unit test caught it immediately. The
round function (below) requires a checkerboard 2-colouring by the parity of $x+y+z$
such that every face-neighbour has the opposite colour. On a torus with wrap-around,
this holds only if the side length is even: an odd-length cycle is not bipartite. With
$N=7$, cells at $x=0$ and $x=6$ are wrap-around neighbours and share a colour,
silently invalidating the Feistel invertibility argument. With $N=8$ the colouring is
consistent, and as a bonus the centre is a literal $2\times2\times2$ cube of 8 cells
--- closer to the original geometric picture than a single point.

\paragraph{Shells.} Since $N$ is even there is no single centre coordinate, only a
centre \emph{pair} $\{3,4\}$ per axis. Define the per-axis shell
$s(v) = \min(|v-3|,|v-4|)$ and the cell shell
$\operatorname{rad}(i) = \max(s(x),s(y),s(z)) \in \{0,1,2,3\}$. Shell~0 is the central
$2\times2\times2$ cube; shell~3 is the outer face shell.

\paragraph{Rate and capacity.} The outer shell ($\operatorname{rad} = 3$, 296 cells) is
the sponge \textbf{rate}: the only cells \texttt{absorb} and \texttt{squeeze} ever
touch. Shells 0--2 (216 cells, 1728 bits) are the \textbf{capacity}: never read or
written directly. This is the legitimate role of ``concentric layers'' promised in
Section~1.

\paragraph{Colouring.} $\operatorname{col}(i) = (x+y+z) \bmod 2$. On the even torus,
every face-neighbour of a cell has the opposite colour.

\subsection{Round function}

One round is $\pi \circ \rho \circ \mathrm{Feistel}$. Each step is individually
invertible, so the composition is a permutation by construction.

\paragraph{$\chi_3$: the nonlinear step.} For cell $i$ with toroidal face-neighbours
$(n_0,\dots,n_5)$ and round $r$, let $j_k = n_{(k + 2(r \bmod 3)) \bmod 6}$. Then
\[
\chi_3(S,i,r) \;=\; (\lnot S_{j_0} \wedge S_{j_1}) \oplus
                   (\lnot S_{j_2} \wedge S_{j_3}) \oplus
                   (\lnot S_{j_4} \wedge S_{j_5}),
\]
computed bytewise. This is in the $\chi$ family of Daemen et al.~\cite{daemen94}: it
has algebraic degree 2, and it is the only source of nonlinearity in the design.
Neighbours are \emph{toroidal}, which gives every cell an identical neighbourhood
shape; without wrapping, shell-boundary cells would diffuse differently from interior
cells, which is precisely the anisotropy the original ``radiate outward'' framing
would have introduced.

\paragraph{Feistel: the invertibility mechanism.} Because every face-neighbour of an
$A$-cell is a $B$-cell, $\chi_3$ evaluated at an $A$-cell reads \emph{only} $B$-cells.
This is exactly the condition a Feistel network needs. With $A_0, B_0$ the colour
classes of the input state:
\begin{align}
  A_1[i] &= A_0[i] \oplus \chi_3(B_0, i, r), &&\text{for all $A$-cells } i,\\
  A_1'   &= A_1 \text{ with } \mathrm{RC}(r) \text{ XORed into cell } 0, \\
  B_1[i] &= B_0[i] \oplus \chi_3(A_1', i, r), &&\text{for all $B$-cells } i.
\end{align}
\begin{proposition}[Invertibility]
The Feistel step is a bijection for \emph{any} function $\chi_3$ whatsoever, provided
$\chi_3$ evaluated at a cell of one colour reads only cells of the other colour.
\end{proposition}
\begin{proof}
Given $(A_1', B_1)$, recover $B_0[i] = B_1[i] \oplus \chi_3(A_1', i, r)$ (which needs
only $A_1'$, available); then $A_1 = A_1'$ with $\mathrm{RC}(r)$ XORed out; then
$A_0[i] = A_1[i] \oplus \chi_3(B_0, i, r)$ (which needs only $B_0$, now recovered).
\end{proof}
This is why we chose a Feistel over a second-order or Margolus partitioning CA: it
decouples the invertibility argument entirely from the choice of nonlinear rule,
letting us make $\chi_3$ as nonlinear as we like without ever endangering decryption.

\paragraph{$\iota$-analogue: round constants.} $\mathrm{RC}(r)$ is a cheap
integer-hash of the round index, XORed into cell 0 only (as Keccak's $\iota$ touches
one lane). It is load-bearing, not decorative: the $\chi_3$ tap rotation has period 3
in $r$, so \emph{without} the constant, round 0 and round 3 are literally the same
function --- a textbook slide-attack foothold. Our \texttt{slide\_test} harness
verifies exactly this (Section~\ref{sec:res-slide}). Note also that any constant state
$S_i = c$ satisfies $\chi_3(S) = 0$, since $\lnot c \wedge c = 0$; without the round
constant, all 256 constant states would be fixed points of the Feistel.

\paragraph{$\rho$: bit-plane coupling.} Rotate each cell's byte left by
$\rho(i) = (x + 3y + 5z) \bmod 8$. Section~\ref{sec:bugs} explains why this step is
the single most important component in the design.

\paragraph{$\pi$: transposition.} A coordinate shear mod 8:
\[
  x' = x + y,\quad y' = y + z,\quad z' = z + x' \pmod 8.
\]
The system is triangular, so the inverse is immediate: $z = z' - x'$, $y = y' - z$,
$x = x' - y$. No lookup table is needed, and bijectivity is verified by unit test
rather than asserted. $\pi$ deliberately moves cells \emph{across} shells; the
rate/capacity split is positional, exactly as Keccak's rate is ``the first $r$ bits''
while $\pi$ shuffles bits through those positions.

\subsection{Mode}

\texttt{Sponge::absorb} XORs up to 296 bytes into the rate cells and applies
$\mathrm{permute} = $ 12 rounds; \texttt{squeeze} reads rate cells and permutes as
needed. \texttt{keystream(key, nonce)} absorbs the key, permutes, absorbs the nonce,
permutes, and squeezes. \emph{There is no authentication tag.} This is a keystream
generator only, and Section~\ref{sec:limitations} says what to do about that (namely:
do not use it).

\section{Two design bugs, and how they were caught}
\label{sec:bugs}

Both bugs below were caught by measurements returning \emph{impossible} numbers, not
by code review. We report them because they are the most transferable content here.

\subsection{The odd-torus colouring}

Discussed in Section~\ref{sec:design}: with $N=7$, the unit test
\texttt{checkerboard\_neighbors\_always\_opposite\_colour} failed instantly. Had we
merely eyeballed the design, the invertibility proof would have been silently false
for wrap-around neighbours.

\subsection{Bit-plane independence: the $6.25\%$ plateau}
\label{sec:planebug}

Our first design had $\chi_3$ and the Feistel but \emph{no} $\rho$. Avalanche
measurement showed the state difference from a single flipped input bit stuck at
$\approx 6.3\%$ and \emph{perfectly flat} through 16 rounds --- not a slow ramp, a
hard ceiling.

The diagnosis is arithmetic. $\chi_3$ uses bitwise \texttt{AND}/\texttt{NOT}/%
\texttt{XOR} on \texttt{u8} cells. Bitwise operations act on all 8 bits in parallel
but \emph{independently}: bit $k$ of the output depends only on bit $k$ of the inputs.
Without $\rho$, the ``$8\times8\times8$ cube of bytes'' was never a single 4096-bit
state at all --- it was \textbf{eight completely independent, non-interacting 512-bit
binary CAs stacked in the same memory}. A flipped bit could never leave its own
bit-plane. Hence the saturation ceiling is
\[
  \frac{512}{2} \Big/ 4096 = 6.25\%,
\]
which is what we measured. Seven-eighths of the state was dead weight \emph{by
construction}.

\begin{remark}
Our initial hypothesis was wrong, and instructively so. We first attributed the
plateau to insufficient long-range transposition. That diagnosis is refuted by
geometry: an 8-torus has diameter 4 per axis, and a Feistel round moves information 2
hops, so $\chi_3$ alone floods the cube in $\approx 3$ rounds. Reach was never the
problem. The tell was that the plateau was \emph{flat and numerically exact}: a
diffusion-reach limitation produces a slow ramp, not a hard ceiling at a suspiciously
round number. \textbf{When a measurement lands on an exact value, compute what that
value would have to mean.}
\end{remark}

$\rho$ fixes it: face-adjacent cells receive different rotation offsets, so bit $k$ of
a cell comes to depend on bit $k - \rho(j)$ of neighbour $j$, and the planes couple.
This is precisely the job Keccak's $\rho$ does, and it is why Keccak has a $\rho$ at
all. A regression test (\texttt{bitplanes\_are\_coupled}) now asserts that a
single-bit difference reaches $\ge 2$ bit-planes within 4 rounds.

\section{Methodology}
\label{sec:method}

All code is Rust (stable, \texttt{--release}), with geometry lookup tables cached
behind \texttt{OnceLock}; the table optimisation is verified behaviour-preserving by
the fact that all 10 unit tests pass unchanged. Measured cost is
$\approx 4.1\,\mu s$ per round per permutation. Four harnesses:

\begin{description}
  \item[\texttt{avalanche}] Strict avalanche ramp from a single flipped bit, averaged
    over 64 trials, reporting total/rate-shell/capacity-shell difference and the
    number of bit-planes reached.
  \item[\texttt{linearity\_test}] Superposition/affinity smoke test: for affine $L$,
    $L(a)\oplus L(b)\oplus L(0)\oplus L(a\oplus b) = 0$ identically.
  \item[\texttt{slide\_test}] Verifies the round constant breaks the period-3
    self-similarity of the tap rotation.
  \item[\texttt{distinguish}] Statistical distinguishers against the sponge's
    fixed-capacity slice: per-bit bias, single-bit differential propagation, and
    random small-weight linear approximations, reported as $z$-scores.
  \item[\texttt{attack}] Cube attack with genuine key recovery: dependency-cone
    measurement, superpoly extraction, BLR linearity testing, GF(2) Gaussian
    elimination, verified against the real key.
  \item[\texttt{degree}] Numeric-mapping degree bounds, soundness cross-check, and
    deterministic zero-sum search.
\end{description}

\subsection{The attacked variant, and why the weakening is sound}
\label{sec:variant}

\texttt{attack} and \texttt{degree} target a \emph{deliberately weakened} variant:
\[
  \text{state} = \text{key}(16\text{B}) \,\|\, \text{nonce}(64\text{B}) \,\|\,
  \text{capacity}(0) \;\to\; \mathrm{permute}_R \;\to\; \text{read rate},
\]
whereas the shipped \texttt{keystream()} places a \emph{full} permutation between the
key absorb and the nonce absorb, so the attacker's nonce never meets the raw key. The
variant is the standard Ascon/Keccak ``init'' model and is strictly easier to attack.

This direction of weakening is what makes the argument sound. An attack that dies at
round $R$ against the weaker variant implies the shipped construction survives
\emph{at least} $R$ rounds: a lower bound on the real thing's security, which is the
useful direction for a defender. Claiming a \emph{strong} result from a weakened
variant would be illegitimate; claiming the conservative one is not.

\subsection{An important non-experiment}

We do not test ``is any output bit of $\mathrm{permute}$ biased'' on the full state.
$\mathrm{permute}$ is a bijection on $\FF_2^{4096}$, so uniform input yields exactly
uniform output by definition; any measured bias is sampling noise. The meaningful
question is the one an attacker faces --- the pushforward through the sponge's
\emph{fixed-capacity slice}, which is not guaranteed uniform. All distinguishers
target that slice.

\section{Two harness bugs, and how they were caught}
\label{sec:harnessbugs}

\subsection{Non-monotonic degree}

Our first \texttt{attack} reported estimated degrees of $0, 1, 0, 4$ for rounds
$1..4$. Degree cannot decrease with rounds under the naive bound, so this was
impossible on its face. Cause: we drew random cubes from 2240 nonce bits and random
output bits. At low rounds the permutation is \emph{local}, so almost every cube sat
entirely outside the output bit's dependency cone. The sums were zero because those
nonce bits do not influence that output bit \emph{at all} --- we were measuring ``did I
hit the cone,'' not degree. Fix: measure the cones first, and draw cubes only from a
target bit's own nonce-influence set.

\subsection{Sampling a sparse target set}

The corrected harness then reported ``no target'' at round 1. Cause: we sampled 24
of 2368 output bits, and at round 1 only $\approx 4\%$ of output bits have \emph{any}
key bit in their cone. Fix: full scan at rounds 1--2 (with the cheaper key-cone test
as a filter before the nonce-cone test), sampling only at rounds $\ge 3$ where cones
have saturated. This is what turned rank 0 into actual recovered key bits.

\section{Results}
\label{sec:results}

All results reproduced on Linux (Ubuntu, container) and Windows 11 (PowerShell 7.6.3)
with different random seeds and secret keys, with agreement to within sampling noise.

\subsection{Correctness}

All 10 unit tests pass: per-round and full-permutation invertibility, $\rho$/$\pi$
round-trips, $\pi$ bijectivity ($|\{\pi(i)\}| = 512$), $\pi^{-1}\circ\pi = \mathrm{id}$,
stream-cipher round-trip, shell partition sizes ($296 + 216 = 512$), checkerboard
neighbour colouring, and bit-plane coupling.

\subsection{Avalanche}

\begin{table}[h]
\centering
\caption{Single-bit avalanche ramp, mean over 64 trials. The bit-plane column is the
diagnostic for the bug of Section~\ref{sec:planebug}.}
\begin{tabular}{rrrrr}
\toprule
Rounds & Total diff & Rate shell & Capacity shell & Bit-planes hit (of 8) \\
\midrule
0 &  0.02\% &  0.03\% &  0.02\% & 1.00 \\
1 &  0.16\% &  0.17\% &  0.15\% & 4.70 \\
2 &  1.08\% &  1.04\% &  1.13\% & 7.84 \\
3 &  5.94\% &  5.94\% &  5.94\% & 8.00 \\
4 & 23.14\% & 23.21\% & 23.05\% & 8.00 \\
5 & 44.55\% & 44.61\% & 44.48\% & 8.00 \\
6 & 49.75\% & 49.93\% & 49.51\% & 8.00 \\
7 & 50.00\% & 49.95\% & 50.07\% & 8.00 \\
\bottomrule
\end{tabular}
\end{table}

Full avalanche is reached at round 6--7, so $\mathrm{ROUNDS} = 12$ provides roughly a
$2\times$ margin on this metric. Note that $5.94\%$ at round 3 is now a mere waypoint;
in the pre-$\rho$ design it was the permanent ceiling. Critically, the
\textbf{capacity shell tracks the rate shell within noise} ($49.51\%$ vs $49.93\%$ at
round 6), with no permanent lag: the concentric-shell rate/capacity split is doing
real work rather than serving as decoration.

\subsection{Affinity smoke test}

Residual weights (out of 4096 bits) over 200 trials per round count, at rounds
$1,2,4,8,12$: minimum $1952$--$1975$, mean $2029.6$--$2050.2$, i.e.\ pinned at
$\approx 50\%$ from round 1 onward. No trial ever produced an exact affine relation.
This is a \emph{necessary, not sufficient} check; it would fail loudly on a
catastrophically linear design and says nothing about algebraic degree.

\subsection{Slide symmetry}
\label{sec:res-slide}

With round constants disabled, rounds 0 and 3 produce identical output on identical
input (\texttt{true}); with constants enabled, they do not (\texttt{false}).
Confirmed: the tap rotation alone yields a period-3 self-similar round function, and
the round constant is exactly what breaks it.

\subsection{Statistical distinguishers}

The strongest signal was single-bit differential propagation (input bit 0 flipped,
20{,}000 trials, worst output bit reported):

\begin{table}[h]
\centering
\caption{Differential propagation against the fixed-capacity slice. $z$ is deviation
from a 50\% flip probability.}
\begin{tabular}{rrrl}
\toprule
Rounds & Worst out-bit & $z$-score & Flip probability \\
\midrule
1--3 & --- & $\pm141.42$ & $\approx 0\%$ or $100\%$ (deterministic) \\
4 & 12 & $-141.42$ & $\approx 0\%$ (deterministic) \\
5 & 1798 & $-105.57$ & $\approx 12.7\%$ \\
6 & 451 & $-18.38$ & $\approx 43.5\%$ \\
7 & 183 & $-3.90$ & $\approx 48.6\%$ \\
8 & 74 & $-3.54$ & $\approx 48.75\%$ \\
\bottomrule
\end{tabular}
\end{table}

Round 5 is trivially distinguishable. Round 6 is a real but small effect --- $z=-18$
is not ``maybe.'' By round 7 it is at the noise floor. This independently corroborates
the avalanche result (full mixing at 6--7) from a different measurement, which is the
kind of internal consistency one wants before trusting either.

\subsection{Cube attack: key recovery}

\begin{table}[h]
\centering
\caption{Dependency cones. Full scan of all 2368 output bits at rounds 1--2; sampling
thereafter.}
\begin{tabular}{rlrrr}
\toprule
Rounds & Scan & Out-bits with key cone & Mean $|$key cone$|$ & Mean $|$nonce cone$|$ \\
\midrule
1 & full   & 253 of 2368 & 3.3   & 2.6 \\
2 & full   & 939 of 2368 & 6.5   & 9.3 \\
3 & sample & 39 of 48    & 24.3  & 98.2 \\
4 & sample & 48 of 48    & 80.5  & 332.6 \\
5 & sample & 48 of 48    & 122.4 & 491.5 \\
\bottomrule
\end{tabular}
\end{table}

\begin{table}[h]
\centering
\caption{Cube attack key recovery against a 128-bit key. Ranges are across four runs
with different secret keys and seeds; cube selection is randomised, so R=2 succeeds in
roughly half of runs. Every ``bits pinned'' figure was verified correct against the
real key.}
\begin{tabular}{rrrrl}
\toprule
Rounds & Linear superpolys & Rank & Keyspace & Key bits pinned \& verified \\
\midrule
1 & 3--9 & 2--5 & $2^{123}$--$2^{126}$ & 1--4 (all correct) \\
2 & 0--2 & 0--2 & $2^{126}$--$2^{128}$ & 0--2 (all correct) \\
3 & 0 & 0 & $2^{128}$ & --- \\
4 & 0 & 0 & $2^{128}$ & --- \\
5 & 0 & 0 & $2^{128}$ & --- \\
\bottomrule
\end{tabular}
\end{table}

Key recovery works reliably at round 1, intermittently at round 2, and never at round
3 or beyond. The cone table explains the mechanism precisely: a $d$-dimensional cube
linearises degree $\approx d+1$, and degree grows $\approx 4^R$. At round 1 the cone
is $\approx 3$ bits and degree $\le 4$, so a $d=3$ cube linearises trivially. By round
3 degree is $\approx 64$, demanding $2^{64}$ work per superpoly. The attack does not
fail because we gave up; it fails because $2^d$ outruns any budget the moment the cone
opens.

\subsection{Provable degree bounds and zero-sums}

Numeric mapping was cross-validated for soundness before use: the bound claims 1336
output bits have degree 0 at round 1 (i.e.\ provably do not depend on \emph{any} nonce
bit). We sampled 120 such predictions and hammered them with random nonces:
\textbf{0 violations}. Had even one moved, every number below would be worthless.

\begin{table}[h]
\centering
\caption{Numeric-mapping degree bounds, all 512 nonce bits active. The bound is sound
(an over-approximation) and beats the naive $4^R$ bound at round 4.}
\begin{tabular}{rrrrr}
\toprule
Rounds & Naive $4^R$ & Proven bound (max) & Median & Out-bits proven degree 0 \\
\midrule
1 & 4   & 4   & 0   & 1336 of 2368 \\
2 & 16  & 16  & 3   & 176 of 2368 \\
3 & 64  & 62  & 19  & 0 \\
4 & 256 & \textbf{151} & 114 & 0 \\
5 & 512 & 512 (sat.) & 512 & 0 \\
\bottomrule
\end{tabular}
\end{table}

\begin{proposition}[Zero-sum]
If output bit $o$ has provable degree $< d$ in a set of $d$ cube variables, then
summing $o$ over all $2^d$ assignments of that cube is $0$ --- for every key, always.
\end{proposition}
\begin{proof}
Every monomial of degree $< d$ omits at least one cube variable; summing over both
values of an omitted variable cancels the monomial in pairs.
\end{proof}

This is a \emph{deterministic} distinguisher: no statistics, no $z$-score. We verified
sampled claims empirically and every predicted sum was indeed zero --- theory and
measurement agreeing.

\begin{table}[h]
\centering
\caption{Cube \emph{shape} sweep. Entries are output bits with a provable zero-sum.
$0$ means the bound was not tight enough to prove one, not that none exists.}
\label{tab:shapes}
\begin{tabular}{lrrrrrr}
\toprule
Shape & $d$ & Cost & $R{=}4$ & $R{=}5$ & $R{=}6$ & $R{=}7$ \\
\midrule
1 cell, all 8 planes    &  8 & $2^{8}$  & 2088 & 0 & 0 & 0 \\
2 cells adjacent        & 16 & $2^{16}$ & 2360 & 0 & 0 & 0 \\
4 cells adjacent        & 32 & $2^{32}$ & 2368 & 0 & 0 & 0 \\
8 cells adjacent        & 64 & $2^{64}$ & 2368 & \textbf{136} & 0 & 0 \\
1 plane across 16 cells & 16 & $2^{16}$ & 924  & 0 & 0 & 0 \\
1 plane across 32 cells & 32 & $2^{32}$ & 1580 & 0 & 0 & 0 \\
2 planes $\times$ 12 cells & 24 & $2^{24}$ & 1335 & 0 & 0 & 0 \\
Spread, stride-37       & 24 & $2^{24}$ & 1230 & 0 & 0 & 0 \\
\bottomrule
\end{tabular}
\end{table}

Cube \emph{shape} matters more than dimension. Compare ``1 plane across 16 cells''
(924 zero-sums) with ``2 cells adjacent'' (2360) --- identical $d = 16$, a
$2.5\times$ difference in reach. The mechanism follows directly from
Section~\ref{sec:planebug}: $\chi_3$ is bitwise, so bit $k$ only meets bit $k$ of its
neighbours, and $\rho$ is the sole coupling between planes. A cube stacked into few
cells therefore grows degree \emph{slower} than one spread across many, and slower
growth means deeper zero-sums. The $d=64$ cube reaches round 5 at $2^{64}$ work:
unverifiable in practice, but $2^{64} \ll 2^{128}$, so it is a legitimate
distinguisher --- proven by propagation rather than measured.

\section{Discussion}

\subsection{The ladder}

\begin{table}[h]
\centering
\caption{Complete cryptanalytic ladder against \texttt{cube3d-sponge} (12 rounds).}
\begin{tabular}{lcl}
\toprule
Result & Reach & Nature \\
\midrule
Key recovery (cube attack)              & 2 & Verified against real key \\
Provable zero-sum, practical ($\le 2^{16}$) & 4 & Deterministic + empirically confirmed \\
Provable zero-sum, $\le 2^{64}$ cube    & 5 & Deterministic, proof-only \\
Statistical differential bias           & 6 & $z = -18$; evidence, not proof \\
\midrule
\textbf{Design}                         & \textbf{12} & $\approx 2\times$ margin over the best distinguisher \\
\bottomrule
\end{tabular}
\end{table}

The four-round gap between key recovery (2) and distinguishing (6) is not an anomaly
--- it is the expected shape. Distinguishers always reach further than key recovery,
which is why ``I found a bias at round 5'' never implies ``the key falls at round 5.''

\subsection{Calibration against Ascon}

The single most useful calibration point is Ascon's own third-party analysis
table~\cite{asconupdate}, which reports a \textbf{zero-sum distinguisher on the full
12/12-round permutation at $2^{55}$}, an 8/12 integral at $2^{46}$, and 5/12
truncated-differential and rectangle distinguishers. Ascon is nonetheless a NIST
standard~\cite{sp800232}.

The lesson is important and cuts against alarmism: \emph{distinguishers on the
permutation do not break the mode}. A zero-sum on the bare permutation says nothing
directly about the AEAD's confidentiality, because the mode never gives the attacker
the freedom the distinguisher assumes. By that yardstick, our round-4/5 zero-sums on a
12-round permutation are unremarkable --- Ascon has one at 12/12. What would be
alarming is a \emph{key recovery} near full rounds, and ours dies at 2.

\subsection{What the geometry was worth}

Three specific findings bear on the original question.

\begin{enumerate}
  \item \textbf{Concentric shells as rate/capacity: validated.} The capacity shell
    tracks the rate shell within noise. The split does real work.
  \item \textbf{Radiating outward as diffusion: refuted before it was built.} Toroidal
    wrapping, not radial structure, is what gives uniform neighbourhoods. Decoupling
    the diffusion geometry from the security boundary was necessary; the original
    proposal fused them.
  \item \textbf{More shells for ``more complexity'': no effect, and mildly harmful.}
    Extra shells add state and cost, add no nonlinearity, and (without wrapping) add
    neighbourhood classes and hence anisotropy.
\end{enumerate}

The deeper point is that \emph{every} security-relevant property in the final design
came from the correction list --- nonlinearity ($\chi_3$), invertibility (Feistel),
symmetry-breaking ($\iota$), plane coupling ($\rho$), transposition ($\pi$) --- and
none from the geometry. The cube is plumbing. Good plumbing, but plumbing.

\section{Limitations}
\label{sec:limitations}

We state these bluntly because the failure mode of this genre is confident
under-qualification.

\begin{enumerate}
  \item \textbf{This is not a secure cipher and must not be used.} There is no
    authentication tag; it is a keystream generator. Use Ascon~\cite{sp800232} or
    ChaCha20-Poly1305.
  \item \textbf{Upper bounds cannot prove security.} Degree upper bounds prove
    distinguishers \emph{exist}; they never prove none does. Every ``0'' in
    Table~\ref{tab:shapes} means ``this bound was not tight enough to prove a
    zero-sum,'' not ``no zero-sum exists.''
  \item \textbf{Numeric mapping is the wrong tool; we used it for tractability.}
    Bit-based division property via MILP~\cite{todo16,xiang16} is strictly
    tighter~\cite{zhang21}, but our state is 4096 bits --- $2.5\times$ Keccak's 1600
    and $12.8\times$ Ascon's 320 --- giving $\approx 4\times10^5$ binary variables and
    $\approx 3\times10^5$ constraints for a 5-round model, with no solver available.
    The gap is not academic: Hu et al.~\cite{hu20} show numeric mapping's Trivium bound
    reaches 79 at 793 rounds where the exact degree at 834 rounds is only 78. A proper
    BDP/monomial-prediction model on real hardware could plausibly push the proven
    reach past 5.
  \item \textbf{Negative attack results are weak.} ``No key bits at round 3'' means
    this attack, at these cube dimensions, with this compute, found nothing. Real cube
    attacks use division-property degree bounds and clever cube
    selection~\cite{todo17cube}, not random draws.
  \item \textbf{No differential/linear trail bounds.} We have no proven bound on the
    maximum expected differential characteristic probability, which is table stakes for
    a real proposal.
  \item \textbf{$\chi_3$, $\rho$, $\pi$ and RC were chosen for plausibility, not
    optimised.} The $\rho$ offsets and the $\pi$ shear were not searched for
    branch-number or diffusion optimality.
  \item \textbf{No third-party analysis.} The design and all attacks share an author.
    This is the weakest point of all.
\end{enumerate}

\section{Conclusion and future work}

We set out to answer whether 3D cellular automata on folded concentric-cube grids
would make a cryptosystem ``way more complex and hence more secure.'' The answer is
that the second clause does not follow from the first, and that the geometric
intuition --- while pointing at genuinely the right shape, since Keccak is a 3D local
nonlinear permutation --- contributes nothing on its own. Security came entirely from
a short, enumerable list of properties, every one of which had to be added
deliberately and then \emph{measured}.

The measurement is the contribution. Four harnesses, cross-validated on two operating
systems, produce a coherent ladder: key recovery to round 2, deterministic zero-sums
to round 4 (practical) or 5 ($2^{64}$), statistical bias to round 6, against a
12-round design. Two design bugs and two harness bugs were caught by numbers that were
\emph{impossible} rather than merely surprising --- a flat plateau at exactly $6.25\%$,
and a non-monotonic degree curve. We would emphasise that pattern above any specific
result: build the falsifier first, and when it returns an exact number, work out what
that number would have to mean.

Future work, in descending order of value: (i) a bit-based division property model
using monomial prediction~\cite{hu20} on hardware capable of a 4096-bit state, to
replace our sampled cones and loose bounds with provable ones; (ii) differential and
linear trail bounds via MILP; (iii) a proper AEAD mode with a tag; (iv) optimisation
of $\rho$/$\pi$ for measured branch number. None of this would make the design
trustworthy without independent cryptanalysis, which remains the only thing that ever
does.

\section*{Availability}

Source, all six harnesses, and the unit-test suite are available at
\url{https://github.com/<user>/3d-ca-crypto}. Results reproduce with
\texttt{cargo test --release} and \texttt{cargo run --release --bin <harness>}.

\begin{thebibliography}{99}

\bibitem{kerckhoffs} A. Kerckhoffs. La cryptographie militaire.
\emph{Journal des sciences militaires}, IX:5--83, 1883.

\bibitem{shannon49} C. E. Shannon. Communication theory of secrecy systems.
\emph{Bell System Technical Journal}, 28(4):656--715, 1949.

\bibitem{wolfram85} S. Wolfram. Cryptography with cellular automata.
In \emph{CRYPTO 1985}, LNCS 218, pp. 429--432. Springer, 1986.

\bibitem{meier91} W. Meier and O. Staffelbach. Analysis of pseudo random sequences
generated by cellular automata. In \emph{EUROCRYPT 1991}, LNCS 547, pp. 186--199.
Springer, 1991.

\bibitem{daemen94} J. Daemen, R. Govaerts and J. Vandewalle. Invertible shift-invariant
transformations on binary arrays. \emph{Applied Mathematics and Computation},
62:259--277, 1994.

\bibitem{camariot} L. Mariot and A. Leporati et al. Insights gained after a decade of
cellular automata-based cryptography. arXiv:2405.02875, 2024.

\bibitem{maiti17} S. Maiti, S. Ghosh and D. Roy Chowdhury. On the security of designing
a cellular automata based stream cipher. In \emph{ACISP 2017}, LNCS 10343. Springer, 2017.

\bibitem{4nbr} Investigating four neighbourhood cellular automata as better cryptographic
primitives. 2017.

\bibitem{fridrich98} J. Fridrich. Symmetric ciphers based on two-dimensional chaotic maps.
\emph{International Journal of Bifurcation and Chaos}, 8(6):1259--1284, 1998.

\bibitem{solak10} E. Solak, C. Çokal, O. T. Yildiz and T. Biyikoğlu. Cryptanalysis of
Fridrich's chaotic image encryption. \emph{International Journal of Bifurcation and
Chaos}, 20(5):1405--1413, 2010.

\bibitem{chen2004} G. Chen, Y. Mao and C. K. Chui. A symmetric image encryption scheme
based on 3D chaotic cat maps. \emph{Chaos, Solitons \& Fractals}, 21(3):749--761, 2004.

\bibitem{lipermonly} C. Li, G. Chen and K.-T. Lo. Cryptanalysis of chaos-based image
encryption algorithms employing different permutation-substitution designs; and
C. Li and K.-T. Lo, Optimal quantitative cryptanalysis of permutation-only multimedia
ciphers against plaintext attacks. \emph{Signal Processing}, 91(4):949--954, 2011.

\bibitem{li09compound} C. Li, S. Li, G. Chen and W. A. Halang. Cryptanalysis of an image
encryption scheme based on a compound chaotic sequence. \emph{Image and Vision Computing},
27(8):1035--1039, 2009.

\bibitem{li11breaking} C. Li, M. Z. Chen and K.-T. Lo. Breaking an image encryption
algorithm based on chaos. \emph{International Journal of Bifurcation and Chaos},
21(7):2067--2076, 2011.

\bibitem{limodified} C. Li, S. Li and K.-T. Lo. Breaking a modified substitution-diffusion
image cipher based on chaotic standard and logistic maps.
\emph{Communications in Nonlinear Science and Numerical Simulation}, 2011.

\bibitem{li18ieaie} C. Li et al. Cryptanalysis of a chaotic image encryption algorithm
based on information entropy. \emph{IEEE Access}, 2018.

\bibitem{rey2015} Á. Martín del Rey. A method to encrypt 3D solid objects based on
three-dimensional cellular automata. In \emph{HAIS 2015}. Springer, 2015.

\bibitem{3dca2016} A study on encryption using three-dimensional cellular automata.
\emph{ScienceAsia}, 42S, 2016.

\bibitem{carev2026} Image encryption with cellular automata: a comprehensive review.
\emph{Archives of Computational Methods in Engineering}, 2026.

\bibitem{spongeindiff} G. Bertoni, J. Daemen, M. Peeters and G. Van Assche. On the
indifferentiability of the sponge construction. In \emph{EUROCRYPT 2008}, LNCS 4965,
pp. 181--197. Springer, 2008.

\bibitem{keccak} G. Bertoni, J. Daemen, M. Peeters and G. Van Assche.
The Keccak sponge function family main document. \url{https://keccak.team}, 2011.

\bibitem{fips202} NIST. FIPS 202: SHA-3 Standard: Permutation-Based Hash and
Extendable-Output Functions. August 2015.

\bibitem{asconspec} C. Dobraunig, M. Eichlseder, F. Mendel and M. Schläffer.
Ascon v1.2: Submission to NIST Lightweight Cryptography. 2021.

\bibitem{asconupdate} C. Dobraunig, M. Eichlseder, F. Mendel and M. Schläffer.
Ascon status update (finalist round). NIST Lightweight Cryptography, 2022.

\bibitem{sp800232} NIST. SP 800-232: Ascon-Based Lightweight Cryptography Standards for
Constrained Devices. 2025.

\bibitem{cube09} I. Dinur and A. Shamir. Cube attacks on tweakable black box polynomials.
In \emph{EUROCRYPT 2009}, LNCS 5479, pp. 278--299. Springer, 2009.

\bibitem{cubetesters} J.-P. Aumasson, I. Dinur, W. Meier and A. Shamir. Cube testers and
key recovery attacks on reduced-round MD6 and Trivium. In \emph{FSE 2009}, LNCS 5665,
pp. 1--22. Springer, 2009.

\bibitem{blr} M. Blum, M. Luby and R. Rubinfeld. Self-testing/correcting with applications
to numerical problems. \emph{Journal of Computer and System Sciences}, 47(3):549--595, 1993.

\bibitem{todo15} Y. Todo. Structural evaluation by generalized integral property.
In \emph{EUROCRYPT 2015}, LNCS 9056, pp. 287--314. Springer, 2015.

\bibitem{todo16} Y. Todo and M. Morii. Bit-based division property and application to
Simon family. In \emph{FSE 2016}, LNCS 9783, pp. 357--377. Springer, 2016.

\bibitem{xiang16} Z. Xiang, W. Zhang, Z. Bao and D. Lin. Applying MILP method to searching
integral distinguishers based on division property for 6 lightweight block ciphers.
In \emph{ASIACRYPT 2016}, LNCS 10031, pp. 648--678. Springer, 2016.

\bibitem{todo17cube} Y. Todo, T. Isobe, Y. Hao and W. Meier. Cube attacks on non-blackbox
polynomials based on division property. In \emph{CRYPTO 2017}, LNCS 10403. Springer, 2017.

\bibitem{liu17} M. Liu. Degree evaluation of NFSR-based cryptosystems.
In \emph{CRYPTO 2017}, LNCS 10403, pp. 227--249. Springer, 2017.

\bibitem{hu20} K. Hu, S. Sun, M. Wang and Q. Wang. An algebraic formulation of the division
property: revisiting degree evaluations, cube attacks, and key-independent sums.
In \emph{ASIACRYPT 2020}, LNCS 12491, pp. 446--476. Springer, 2020.

\bibitem{zhang21} S. Zhang, X. Zeng et al. On the relationships between different methods
for degree evaluation. \emph{IACR Transactions on Symmetric Cryptology}, 2021(1).

\bibitem{boura11} C. Boura, A. Canteaut and C. De Cannière. Higher-order differential
properties of Keccak and Luffa. In \emph{FSE 2011}, LNCS 6733, pp. 252--269. Springer, 2011.

\bibitem{canteaut02} A. Canteaut and M. Videau. Degree of composition of highly nonlinear
functions and applications to higher order differential cryptanalysis.
In \emph{EUROCRYPT 2002}, LNCS 2332, pp. 518--533. Springer, 2002.

\end{thebibliography}

\end{document}
